\begin{frame}{Aktueller Stand: Eventloop- und Netzwerkschicht}
  \begin{itemize}
    \item \alert{Unterst\"utzte Vermittlungsprotokolle} (OSI Schicht 3)
    \begin{itemize}
      \item IPv4
      \item IPv6
      \item Bluetooth l2cap (Logical Link Control Adaptation Protokoll)
      \item Bluetooth rfcomm (Radio Frequency Communication)
    \end{itemize}
    \item \alert{Transportprotokolle} (OSI Schicht 4)
    \begin{itemize}
      \item \alert{Unterst\"utzt}
      \begin{itemize}
        \item \alert{SOCK\_STREAM} (bedeutet TCP im Falle von IPv4 und IPv6)
      \end{itemize}
      \item \alert{Noch nicht unterst\"utzt}
      \begin{itemize}
        \item \alert{SOCK\_DGRAM} (bedeutet UDP im Falle von IPv4 und IPv6)
      \end{itemize}
    \end{itemize}
    \item \alert{Security}
    \begin{itemize}
      \item \alert{SSL/TLS} mit server- und clientseitiger Zertifikatspr\"ufung (X.509).
      \begin{itemize}
        \item Direkt im Transportlayer implementiert.
        \item Verwendbar auf jedem Vermittlungsprotokoll.
      \end{itemize}
    \end{itemize}
    \item \alert{Timer}
    \begin{itemize}
      \item Singleshot- und Intervalltimer (API wie in node.js)
    \end{itemize}
    \item {\alert{API} für darüber liegende Schichten}
    \begin{itemize}
      \item \alert{Server}-Klassen (als C++ Template)
      \item \alert{Client}-Klassen (als C++ Template)
    \end{itemize}
  \end{itemize}
\end{frame}
